%% ma: L10-B13-0102
%% lop: 10
%% chuong: 5
%% bai: 13
%% dang: 10.13.13
%% loai: DS
%% muc: [NB, TH, TH, TH]
%% dap_an: "ĐSSĐ"
%% nguon: "chuyen-de-cac-so-dac-trung-cua-mau-so-lieu-khong-ghep-nhom-toan-10-chuong-trinh-moi.pdf (GV. Phan Nhật Linh) (P03), Bài 02 Dạng 2 Phần II Câu 6"
%% kien_thuc: "Lớp 10 Bài 13 (số trung bình, trung vị, tứ phân vị, mốt)"
%% trang_thai: chinh-thuc
%% ghi_chu: "Tự duyệt theo yêu cầu giáo viên 22/10/2026."
\begin{ex}
Cho bảng số liệu thống kê chiều cao (cm) của một nhóm học sinh như sau:
\begin{center}
\begin{tabular}{|c|c|c|c|c|c|c|c|}\hline
$150$ & $153$ & $153$ & $154$ & $154$ & $155$ & $160$ & $160$\\\hline
$162$ & $162$ & $163$ & $163$ & $163$ & $165$ & $165$ & $167$\\\hline
\end{tabular}
\end{center}
Xét tính đúng sai của các khẳng định sau:
\choiceTF{\True Có $16$ học sinh được thống kê chiều cao}{Chiều cao trung bình của nhóm học sinh là $160$ cm}{Trung vị của mẫu số liệu là $160$}{\True Tứ phân vị thứ nhất của mẫu số liệu là $154$}
\loigiai{a) Đúng.

b) Sai: $\overline x=\dfrac{2549}{16}=159{,}3125$.

c) Sai: $M_e=\dfrac{x_8+x_9}2=\dfrac{160+162}2=161$.

d) Đúng: $Q_1=\dfrac{x_4+x_5}2=\dfrac{154+154}2=154$.}
\end{ex}
